\chapter{Marchés physiques de matières premières}\label{ch:m3:physical-commodity-markets}

Un pétrolier termine le chargement de 700\,000 barils de brut de la mer du Nord un mardi.
L'acheteur et le vendeur ont signé leur contrat il y a plusieurs semaines, et ni l'un ni l'autre
ne connaît encore le prix. Le contrat stipule~: la moyenne d'une évaluation quotidienne publiée
du brut de la mer du Nord sur cinq jours ouvrés, deux jours avant la date figurant sur le
document de chargement, cette date, et deux jours après, plus quarante cents. Les deux parties
connaîtront le prix le lundi soir suivant, lorsque la dernière des cinq
évaluations sera publiée. D'ici là, la raffinerie qui a acheté le pétrole détient une cargaison dont le coût
varie chaque jour, et elle doit décider si, et comment, la fixer. Ce chapitre décrit la couche
du monde des matières premières que forment de tels contrats, le \omterm{def:m3:physical-commodity-markets:layers}{marché physique}, et les outils
grâce auxquels les traders transforment ses prix flottants en prix fixes sur lesquels une
entreprise peut bâtir ses prévisions.

\section{Physique et papier}

Les marchés de matières premières comptent deux couches. L'une concerne la circulation des biens~: navires, pipelines, citernes,
silos et entrepôts, avec des contrats qui précisent une qualité, une quantité, un lieu
et une date. L'autre ne concerne que l'argent~: contrats à terme, swaps et options dont
le règlement dépend d'un prix. Les livres~1 et~2 de cette série ont traité presque
exclusivement du second type d'instrument~; pour les matières premières, le premier type passe au premier plan,
parce que c'est là que les prix se forment.

\begin{definition}[Marché physique, marché papier]\label{def:m3:physical-commodity-markets:layers}
Le \emph{marché physique}\index{marché physique} (physical market) d'une matière première est l'ensemble des
contrats de livraison d'une quantité et d'une qualité déterminées de la matière première en
un lieu et à un moment donnés. Son \emph{marché papier}\index{marché papier} (paper market) est l'ensemble
des contrats portant sur le prix de la matière première, qui sont normalement réglés en espèces ou clôturés
avant la livraison~: contrats à terme, swaps et options.
\end{definition}

La frontière est poreuse. Un contrat à terme avec livraison physique
(livre~1, chapitre~21) est un papier pendant presque toute sa durée de vie et devient physique
à son échéance~; un exchange for physical fait passer de l'un à l'autre. Le \omterm{def:m3:physical-commodity-markets:layers}{marché papier}
a pour outil de travail, en dehors des bourses, le swap.

\begin{definition}[Swap sur matières premières]\label{def:m3:physical-commodity-markets:swap}
Un \emph{swap sur matières premières}\index{swap sur matières premières} (commodity swap) échange, pour chaque période de sa
durée de vie, un prix fixe contre la moyenne d'un prix flottant publié sur la période,
appliqués l'un et l'autre à un volume convenu, et se règle en espèces sur la différence.
\end{definition}

\begin{example}[Un swap mensuel sur le brut]\label{ex:m3:physical-commodity-markets:swap}
Une compagnie aérienne achète un swap portant sur 100\,000 barils par mois pendant le trimestre suivant, à un prix fixe de
\qty{82.00}{\$/\barrel} contre la moyenne mensuelle d'un indice de référence du brut. Si
l'indice affiche en moyenne \qty{85.40}{\$/\barrel} le premier mois, la compagnie
reçoit $3.40 \times 100\,000 = \$340\,000$~; si la moyenne est de \qty{79.10}{\$/\barrel}
le deuxième mois, elle verse $\$290\,000$. Quel que soit le prix qu'elle paie pour son carburant sur le
\omterm{def:m3:physical-commodity-markets:layers}{marché physique}, son coût net par baril est fixé autour de 82~\$, à la différence près
entre le carburant qu'elle brûle et l'indice de référence qu'elle a couvert, ce qui fait l'objet de
la dernière section.
\end{example}

\section{Maisons de négoce et chaîne de valeur physique}

Entre le producteur qui extrait une matière première et le consommateur qui la brûle, la fond
ou la moud, se trouvent des entreprises qui achètent, stockent, mélangent, expédient et revendent.

\begin{definition}[Maison de négoce]\label{def:m3:physical-commodity-markets:house}
Une \emph{maison de négoce}\index{maison de négoce} (commodity trading house) est une entreprise dont
l'activité consiste à acheter des matières premières physiques et à les revendre, en gagnant la différence
entre leur valeur en un lieu, à un moment ou sous une forme donnée et leur valeur ailleurs, et
à couvrir le prix nu par du papier.
\end{definition}

Une \omterm{def:m3:physical-commodity-markets:house}{maison de négoce} transforme une matière première de trois façons, chacune rémunérée
par un écart de prix~: dans l'espace (une cargaison achetée à un port de chargement et vendue à un
port de déchargement), dans le temps (du pétrole stocké lorsque le prix à terme rémunère la citerne)
et sous forme (des qualités mélangées pour obtenir une spécification que veut une raffinerie). L'étude de Pirrong
sur le secteur décrit les sociétés de trading exactement dans ces termes, comme
des acteurs de transformations physiques, dont les profits proviennent des spreads
entre lieux, dates et qualités plutôt que du niveau des prix.
C'est pourquoi une \omterm{def:m3:physical-commodity-markets:house}{maison de négoce} bien gérée couvre presque toute son exposition au prix nu~:
le niveau des cours pétroliers est un risque qu'elle porte, et non celui qui
la rémunère.

\begin{dated}{2026-09}{Quelle est la taille du plus grand négociant}\label{dat:m3:physical-commodity-markets:scale}
Vitol, parmi les plus grands négociants en énergie, a annoncé des livraisons de 8 millions de barils par jour
de brut et de produits en 2025 (7.2 millions en 2024), pour un chiffre d'affaires de
343~milliards de dollars USD. La demande mondiale de pétrole s'élevait à environ 104 millions de barils par jour~; l'Agence internationale
de l'énergie a prévu en mai 2026 qu'elle reculerait en 2026, pour atteindre 104 millions
de barils par jour, après que les expéditions par le détroit d'Ormuz ont été perturbées.
\end{dated}

\begin{remark}[Pourquoi l'activité est concentrée]\label{rem:m3:physical-commodity-markets:scale}
Un négoce physique immobilise un fonds de roulement pendant toute la traversée~: une
cargaison de 700\,000 barils à \qty{80}{\$/\barrel} représente 56~millions de dollars USD, payés à
la livraison et récupérés quelques semaines plus tard. Les entreprises qui le font à grande échelle sont celles
qui disposent des lignes de crédit bancaire pour le financer (\cref{ch:m3:getting-access-commodities-and-power})
et du flux nécessaire pour répartir leurs coûts fixes sur de nombreuses cargaisons.
\end{remark}

\section{Conditions de livraison et contrat physique}

Un contrat physique doit préciser où la marchandise change de mains, qui paie son transport
et qui supporte le risque de perte pendant le trajet. La pratique internationale
recourt à des conditions types pour cela.

\begin{definition}[Incoterms, franco à bord, coût, assurance et fret]\label{def:m3:physical-commodity-markets:incoterms}
Les \emph{Incoterms}\index{Incoterms} sont les conditions de livraison types publiées par la
Chambre de commerce internationale. Dans le cadre du \emph{franco à bord}\index{franco à bord}
(FOB), le vendeur livre la marchandise à bord du navire désigné par l'acheteur,
au port d'expédition convenu~; le risque est transféré à cet endroit, et l'acheteur paie le fret
et l'assurance. Dans le cadre du \emph{coût, assurance et fret}\index{coût, assurance et fret}
(CIF), le risque est également transféré à bord au port d'expédition, mais le vendeur
conclut et paie le transport jusqu'au port de destination convenu, ainsi que l'assurance
contre le risque de perte supporté par l'acheteur.
\end{definition}

\begin{definition}[Connaissement]\label{def:m3:physical-commodity-markets:bl}
Un \emph{connaissement}\index{connaissement} (bill of lading) est le document que le transporteur émet
au chargement~: un reçu de la marchandise, une preuve du contrat de transport et un titre
de propriété dont le porteur peut réclamer la cargaison à destination. Sa date, la date du
B/L, sert d'ancrage à la plupart des formules de prix.
\end{definition}

\begin{omfigure}
\centering
\begin{tikzpicture}[font=\small,
  port/.style={draw,thick,rounded corners=2pt,align=center,minimum height=1.0cm,minimum width=2.6cm},
  arr/.style={-{Stealth[length=2.5mm]},thick}]
\node[port,draw=omDef,fill=omDef!8] (load) at (0,0) {port de chargement\\(citerne du vendeur)};
\node[port,draw=omThm,fill=omThm!8] (disch) at (9.4,0) {port de déchargement\\(raffinerie de l'acheteur)};
\draw[arr] (load.east) -- node[above=2pt] {traversée~: fret, assurance, pertes} (disch.west);
\draw[omProp,very thick] (2.05,-0.9) -- (2.05,0.9);
\node[omProp,font=\footnotesize\bfseries,align=center] at (2.05,1.25) {transfert du risque à bord\\(FOB et CIF)};
\node[anchor=west,font=\footnotesize] at (-1.3,-1.35) {\textbf{FOB~:} le vendeur paie jusqu'au bastingage du navire~; l'acheteur affrète, assure et paie le fret};
\node[anchor=west,font=\footnotesize] at (-1.3,-1.85) {\textbf{CIF~:} le vendeur paie aussi le fret et l'assurance jusqu'à destination~; le risque passe néanmoins à bord};
\node[anchor=west,font=\footnotesize] at (-1.3,-2.35) {\textbf{Date du B/L~:} le jour où le connaissement est émis~; les périodes de prix se comptent à partir de celle-ci};
\end{tikzpicture}

\omcaption{Les deux conditions de livraison du commerce maritime de matières premières. Dans les deux cas,
l'acheteur supporte le risque du transport~; sous CIF, le vendeur en paie le coût et le répercute
dans le prix de la cargaison. Schéma simplifié.}\label{fig:m3:physical-commodity-markets:incoterms}
\end{omfigure}

Un prix CIF est donc approximativement égal au prix FOB augmenté du fret et de l'assurance~:
les deux conditions diffèrent par la partie qui organise et paie le transport, et non par
celle qui en supporte le risque. Au-delà de l'Incoterm, un contrat physique désigne la qualité et sa
spécification, le volume assorti d'une tolérance opérationnelle (une cargaison est rarement
exactement la taille convenue), la plage de dates de chargement ou de livraison,
l'agent d'inspection dont le mesurage de la quantité et de la qualité fait foi, les modalités de paiement
(souvent par lettre de crédit, trente jours après la date du B/L~; \cref{ch:m3:getting-access-commodities-and-power}) et la
formule de prix.

\section{Agences de cotation des prix et fenêtres d'évaluation}

La plupart des cargaisons physiques ne se négocient pas en bourse, si bien qu'il n'existe pas
de prix de bourse à inscrire dans leurs formules. Les prix proviennent d'éditeurs spécialisés.

\begin{definition}[Agence de cotation des prix, fenêtre d'évaluation]\label{def:m3:physical-commodity-markets:pra}
Une \emph{agence de cotation des prix}\index{agence de cotation des prix} (PRA) est un éditeur
qui évalue et publie, chaque jour ouvré, les prix des matières premières physiques
et des contrats de gré à gré qui leur sont liés, à partir des cours acheteurs, des cours vendeurs et des transactions qu'elle observe
et des déclarations des intervenants du marché. Une \emph{fenêtre
d'évaluation}\index{fenêtre d'évaluation} (assessment window) est une période quotidienne fixe pendant laquelle l'agence
recueille les cours acheteurs, les cours vendeurs et les transactions qu'elle utilise ensuite pour fixer son évaluation à
l'heure de clôture de la fenêtre.
\end{definition}

Les deux plus grandes agences dans le pétrole sont Platts, filiale de S\&P Global, et Argus. Leurs
évaluations sont citées dans les contrats physiques et par les contrats papier qui
les couvrent, et servent donc d'indices de référence. En octobre 2012, à la demande du
G20, l'Organisation internationale des commissions de valeurs a publié
des principes applicables aux \omterm{def:m3:physical-commodity-markets:pra}{agences de cotation des prix} du pétrole dont les évaluations sont citées par
des contrats dérivés~: la gouvernance, la qualité et l'intégrité de la méthodologie,
et la gestion des conflits d'intérêts.

\begin{dated}{2026-09}{La fenêtre de la mer du Nord}\label{dat:m3:physical-commodity-markets:moc}
Platts évalue le Dated Brent à partir de son processus de clôture de marché~: les cours acheteurs, les cours vendeurs et les transactions déclarés pendant une fenêtre de trente minutes,
de 16~h~00 à 16~h~30, heure de Londres, chaque jour ouvré,
avec une évaluation qui reflète la valeur de la qualité la plus compétitive à 16~h~30.
Les participants soumettent des cours acheteurs et des cours vendeurs pour des cargaisons de taille donnée, de dates de chargement et
de qualité déterminée~; l'agence publie ces cours et les transactions. Une fenêtre parallèle fonctionne à Singapour
pour le brut du Moyen-Orient.
\end{dated}

\section{Formules de prix, netbacks et couverture d'une opération physique}

\begin{definition}[Période de prix, différentiel de prix]\label{def:m3:physical-commodity-markets:formula}
La \emph{période de prix}\index{période de prix} (pricing period) d'un contrat physique est l'ensemble des
jours dont les évaluations de référence sont moyennées pour fixer son prix~: par exemple les cinq
jours ouvrés autour de la date du B/L (notée 2-1-2), ou tous les jours du
mois de livraison. Le \emph{différentiel de prix}\index{différentiel de prix} (price differential) est la prime fixe
ou la décote ajoutée à cette moyenne~; elle traduit la valeur de la qualité, du lieu et
du calendrier par rapport à l'indice de référence, et c'est ce que les parties négocient réellement.
\end{definition}

Un trader qui achète en un lieu et vend en un autre les compare grâce à leur
valeur en un point unique.

\begin{definition}[Netback]\label{def:m3:physical-commodity-markets:netback}
Le \emph{netback}\index{netback} d'une vente est son prix livré diminué des coûts de
l'acheminement de la marchandise jusqu'à destination (fret, assurance, pertes en cours de transport, droits et
frais de terminal), exprimé par unité au point d'origine.
\end{definition}

\begin{example}[Choisir une destination par netback]\label{ex:m3:physical-commodity-markets:netback}
Une cargaison peut être vendue livrée à Rotterdam à \qty{81.30}{\$/\barrel} avec
\qty{1.10}{\$/\barrel} de fret, ou à Singapour à \qty{83.10}{\$/\barrel} avec
\qty{3.05}{\$/\barrel}~; l'assurance coûte \qty{0.03}{\$/\barrel} dans les deux cas. Les
\omterm{def:m3:physical-commodity-markets:netback}{netbacks} s'élèvent à \qty{80.17}{\$/\barrel} et \qty{80.02}{\$/\barrel}~: Rotterdam l'emporte, bien
que Singapour paie davantage.
\end{example}

Une formule physique laisse les parties exposées à l'indice de référence jusqu'au
dernier jour de fixation. Elles se couvrent contre cette exposition avec des contrats papier sur le même indice de référence, ou sur un
indice très proche, et ce que la couverture ne peut pas supprimer a un nom.

\begin{definition}[Risque de base]\label{def:m3:physical-commodity-markets:basisrisk}
Le \emph{risque de base}\index{risque de base} (basis risk) est le risque que le prix de la position
couverte et le prix de l'instrument de couverture évoluent de montants différents~: le
risque qui subsiste dans une position couverte parce que la base entre les deux change.
\end{definition}

\begin{method}[Couverture d'une cargaison dont le prix est fixé sur une période]\label{met:m3:physical-commodity-markets:hedge}
Lorsque le prix est fixé sur $n$ jours de fixation, chaque jour fixe $1/n$ du volume.
(i) Dès la création de l'exposition (le contrat est signé et son prix va
flotter), prendre une position sur contrats à terme de même taille et de signe opposé à
l'exposition que la fixation va créer~: un acheteur de pétrole à prix flottant qui veut un coût fixe achète
des contrats à terme dès maintenant. (ii) À chaque jour de fixation, clôturer $1/n$ des contrats à terme
au règlement-livraison de ce jour. (iii) Le prix couvert est alors égal au prix à terme initial plus la
base moyenne sur les jours de fixation, plus le différentiel.
\end{method}

\begin{proposition}[Ce que coûte une cargaison couverte]\label{prop:m3:physical-commodity-markets:hedged}
Soient $F_d$ le prix à terme et $b_d$ la base, égale au prix physique moins le prix à terme, le
jour $d$, et soient $d_1,\dots,d_n$ les jours de fixation. Un acheteur couvert selon
la \cref{met:m3:physical-commodity-markets:hedge} à partir du jour $0$ paie par baril
\[
\frac1n\sum_{i=1}^n (F_{d_i}+b_{d_i}) + D - \frac1n\sum_{i=1}^n (F_{d_i}-F_0)
\;=\; F_0 + \frac1n\sum_{i=1}^n b_{d_i} + D ,
\]
où $D$ est le différentiel. Le prix flat a disparu~; seule subsiste la base moyenne.
\end{proposition}

\begin{proof}
Le prix par formule est la première moyenne plus $D$~; les contrats à terme achetés à $F_0$ et
vendus par parts égales en $F_{d_i}$ rapportent la seconde moyenne. Il reste à soustraire.
\end{proof}

\Cref{fig:m3:physical-commodity-markets:costcdf} simule la raffinerie du
hook~: les contrats à terme évoluent de \qty{1.80}{\$/\barrel} par jour et la base de
\qty{0.12}{\$/\barrel} par jour, toutes deux sous forme de marches aléatoires. Sans couverture, le coût de la
cargaison présente un écart type de \qty{8.13}{\$/\barrel}, soit environ 5.7~millions de dollars sur
la cargaison~; couverte, elle est de \qty{0.54}{\$/\barrel}, soit exactement l'écart type
de la base moyenne sur les cinq jours de fixation.

\begin{omfigure}
\begin{tikzpicture}
\begin{axis}[width=11cm,height=5cm,
  xlabel={coût de la cargaison, \$/bbl},ylabel={probabilité cumulée},
  tick label style={font=\small},label style={font=\small},
  xmin=58,xmax=102,ymin=0,ymax=1,grid=major,grid style={gray!25},
  legend columns=2,legend style={font=\footnotesize,at={(0.5,-0.3)},anchor=north,draw=none,fill=none,/tikz/every even column/.append style={column sep=8pt}},
  table/col sep=comma]
\addplot[omDef,very thick] table[x=unhedged,y=p]{figdata/markets-3/01-physical-commodity-markets/costcdf.csv};
\addplot[omThm,very thick] table[x=hedged,y=p]{figdata/markets-3/01-physical-commodity-markets/costcdf.csv};
\legend{non couvert,couvert avec contrats à terme sur la période de prix}
\end{axis}
\end{tikzpicture}

\omcaption{Distribution du coût d'une cargaison fixée en 2-1-2 autour d'une date de \omterm{def:m3:physical-commodity-markets:bl}{connaissement} à vingt
jours ouvrés, avec contrats à terme à \qty{80}{\$/\barrel}, une base de
$-\qty{0.30}{\$/\barrel}$ et un différentiel de $+\qty{0.40}{\$/\barrel}$~: 20\,000
trajectoires simulées. La couverture élimine le prix flat et laisse subsister la base.
Volatilités illustratives~; données~: le tutoriel du chapitre.}\label{fig:m3:physical-commodity-markets:costcdf}
\end{omfigure}

Le \omterm{def:m3:physical-commodity-markets:basisrisk}{risque de base} est faible lorsque la couverture et la cargaison se réfèrent au même indice de référence et
élevé lorsque ce n'est pas le cas. Les principales références physiques, le baril maritime de la mer du Nord
et le baril continental américain (\cref{ch:m3:crude-oil}), évoluent ensemble la plupart du
temps et s'écartent lorsque la logistique se grippe
(\cref{fig:m3:physical-commodity-markets:differential}). En 2011, la production américaine de brut continental
dépassa les capacités des oléoducs sortant du hub de stockage de Cushing, dans l'Oklahoma, et
le baril continental est tombé à une décote mensuelle moyenne d'environ 27~\$ en septembre
2011~; en 2026, le baril maritime est monté au-dessus de lui de 17.66~\$ au plus dans la
moyenne mensuelle d'avril, pendant que les expéditions par le détroit d'Ormuz étaient perturbées.
Un hedger ayant vendu un indice de référence contre une cargaison fixée sur l'autre aurait
perdu ou gagné ces montants en plus du prix flat.

\begin{omfigure}
\begin{tikzpicture}
\begin{axis}[width=11cm,height=4.8cm,
  ylabel={WTI moins Brent, \$/bbl},
  tick label style={font=\small},label style={font=\small},
  xmin=2000,xmax=2026.75,ymin=-32,ymax=10,ytick={-30,-20,-10,0,10},grid=major,grid style={gray!25},
  xtick={2000,2004,2008,2012,2016,2020,2024},xticklabel style={/pgf/number format/1000 sep={}},
  table/col sep=comma]
\addplot[omDef,thick] table[x=t,y=spread]{figdata/markets-3/01-physical-commodity-markets/differential.csv};
\draw[gray] (axis cs:2000,0) -- (axis cs:2026.75,0);
\end{axis}
\end{tikzpicture}

\omcaption{L'indice de référence américain continental moins l'indice de référence maritime de la mer du Nord, moyennes mensuelles
des prix spot quotidiens, de janvier 2000 à août 2026. L'engorgement des oléoducs de
2011 et la perturbation du Golfe de 2026 sont les deux décotes les plus profondes. Données~: séries FRED
DCOILWTICO et DCOILBRENTEU (US Energy Information Administration).}\label{fig:m3:physical-commodity-markets:differential}
\end{omfigure}

\section{Tutoriel~: fixation du prix et couverture d'une cargaison}\label{tut:m3:physical-commodity-markets:hedge}

\begin{tutorial}
\textbf{Objectif.} Fixer le prix d'une cargaison à partir de sa formule, la couvrir par
la \cref{met:m3:physical-commodity-markets:hedge}, et mesurer ce que la couverture laisse subsister.
\textbf{État final~:} \cref{fig:m3:physical-commodity-markets:costcdf} et les
écarts types \qty{8.13}{} et \qty{0.54}{\$/\barrel}.
\begin{tutsteps}
\item \textbf{Les jambes et la couverture.} Une jambe est achetée ou vendue, fixée sur un ensemble de
jours~; chaque jour de fixation fixe une part égale et la couverture passe des ordres pour cette part.
\omcode{code/firm/physdeal/firm_physdeal.py}{28}{58}{Jambes, exposition et transactions sur contrats à terme qui la compensent.}\label{lst:m3:physical-commodity-markets:legs}
\item \textbf{Simuler.} Contrats à terme et base sous forme de marches aléatoires~; coût non couvert et couvert,
et l'écart type théorique de la cargaison couverte, issu de la covariance d'une marche aléatoire.
\omcode{code/markets-3/01-physical-commodity-markets/python/m3_physical.py}{22}{39}{La cargaison de la raffinerie sur 20\,000 trajectoires.}\label{lst:m3:physical-commodity-markets:sim}
\item \textbf{Exécuter} \texttt{m3\_physical.simulate\_cargo()},
\texttt{hedged\_sd\_theory()} et \texttt{fig\_physical.py}.
\end{tutsteps}
\textbf{Pour aller plus loin.} Couvrir à la seule date du \omterm{def:m3:physical-commodity-markets:bl}{connaissement}, au lieu de
sur la \omterm{def:m3:physical-commodity-markets:formula}{période de prix}, et mesurer le risque résiduel~; remplacer la marche aléatoire de la base par une
à retour vers la moyenne et voir quelle part du risque couvert subsiste.
\end{tutorial}

\section{Construction~: le modèle des transactions physiques}\label{bld:m3:physical-commodity-markets:physdeal}

\begin{build}
\bfield{Rôle} Chaque transaction physique de la société miniature entre dans ses livres
par ce modèle~: quel sera le prix, quels jours restent flottants, quel fret et quelle assurance
la société paie, et quels contrats à terme négocier chaque jour pour garder le livre à plat.
\bfield{Interface} \texttt{pricing\_days(event\_day, before, after)}~;
\texttt{formula\_price(assessments, days, differential)}~; \texttt{Leg(side, volume,
days, differential)}~; \texttt{exposure(legs, day)}~; \texttt{hedge\_trades(legs)}~;
\texttt{cost\_split(incoterm, freight, insurance)}~; \texttt{netback}~;
\texttt{locked\_margin}.
\bfield{Règles} Les jours sont des indices de jours ouvrés~; une évaluation manquante est une erreur, jamais
une interpolation~; chaque jour de fixation fixe une part égale~; les lots de contrats à terme sont
(1\,000 barils) ou le calendrier est rejeté~; FOB et CIF uniquement.
\bfield{Tests d'acceptation} \texttt{code/firm/physdeal/tests/}~: une période 2-1-2~;
l'exposition d'une cargaison back-to-back est nulle avant et après les deux \omterm{def:m3:physical-commodity-markets:formula}{périodes de prix} et
le volume plein entre elles~; la somme de la couverture et de l'exposition est nulle chaque jour~; le \omterm{def:m3:physical-commodity-markets:netback}{netback} et
la répartition selon l'Incoterm.
\bfield{Pour aller plus loin} Enregistrer les transactions sur contrats à terme dans le gestionnaire de positions du livre~1,
chapitre~7~; calendriers des dates de prix avec jours fériés~; tolérance de volume au
gré de l'acheteur.
\end{build}

\begin{omsources}
\item International Chamber of Commerce, \emph{Incoterms 2020}~; ICC Academy, «~Place
of delivery and risk transfer in international trade contracts~».
\item IOSCO, \emph{Principles for Oil Price Reporting Agencies}, FR06/12, 5 octobre 2012
(copie sur Internet Archive), et le rapport de mise en œuvre de 2013.
\item S\&P Global Platts, méthodologie de Dated Brent et processus market-on-close.
\item C.~Pirrong, \emph{The Economics of Commodity Trading Firms}, 2014.
\item Vitol, \emph{2025 volumes and review}~; IEA, \emph{Oil Market Report}, mai et
août 2026.
\item US Energy Information Administration, \emph{Today in Energy}, «~Spread narrows
between Brent and WTI crude oil benchmark prices~», 2013~; séries FRED DCOILWTICO et
DCOILBRENTEU.
\end{omsources}

\section{Exercices}

\begin{exercise}[$\star$]\label{exo:m3:physical-commodity-markets:1}
Un cargo est valorisé à la moyenne de cinq évaluations, plus \qty{0.35}{\$/\barrel}.
Les évaluations sont 81.20, 80.70, 79.90, 80.40 et 81.05. Quel est le prix~?
\end{exercise}

\begin{exercise}[$\star$]\label{exo:m3:physical-commodity-markets:2}
Un vendeur propose un cargo à \qty{78.40}{\$/\barrel} FOB. Le fret jusqu'au port de l'acheteur
est de \qty{1.35}{\$/\barrel} et l'assurance de \qty{0.04}{\$/\barrel}. Quel prix CIF est
équivalent, et qui supporte le risque du voyage dans chacun des deux termes~?
\end{exercise}

\begin{exercise}[$\star$]\label{exo:m3:physical-commodity-markets:3}
Un cargo se vend livré à \qty{84.10}{\$/\barrel}~; le fret est de
\qty{2.20}{\$/\barrel}, l'assurance de \qty{0.05}{\$/\barrel}, et 0.2~\% du volume
est perdu en transit. Donner le \omterm{def:m3:physical-commodity-markets:netback}{netback} par baril chargé.
\end{exercise}

\begin{exercise}[$\star\star$]\label{exo:m3:physical-commodity-markets:4}
Une raffinerie achète 700\,000 barils valorisés en 2-1-2 autour de la date du B/L. Combien de
contrats à terme de 1\,000 barils achète-t-elle aujourd'hui, et combien en vend-elle chaque
jour de fixation~? Que change-t-il si l'un des cinq jours est un jour férié où aucune évaluation
n'est publiée et que le contrat l'écarte~?
\end{exercise}

\begin{exercise}[$\star\star$]\label{exo:m3:physical-commodity-markets:5}
Dans la simulation de la \cref{fig:m3:physical-commodity-markets:costcdf}, l'écart-type couvert
est de \qty{0.54}{\$/\barrel}. Quel serait-il si la volatilité quotidienne de la base doublait,
et pourquoi la volatilité des contrats à terme n'intervient-elle pas~?
\end{exercise}

\begin{exercise}[$\star\star$]\label{exo:m3:physical-commodity-markets:6}
En reprenant l'\cref{ex:m3:physical-commodity-markets:netback}, de combien le prix livré à
Singapour devrait-il augmenter pour que le trader change de destination~?
\end{exercise}

\begin{exercise}[$\star\star\star$]\label{exo:m3:physical-commodity-markets:7}
\emph{Programmation.} Avec \texttt{load\_differential}, trouver le mois où la décote du WTI
par rapport au Brent a été la plus forte depuis 2000 et son montant, puis compter les mois depuis
janvier 2011 où le WTI a été en moyenne au-dessus du Brent.
\end{exercise}

\begin{exercise}[$\star\star\star$]\label{exo:m3:physical-commodity-markets:8}
\emph{Repérer l'erreur.} «~Notre cargo est valorisé autour de la date du B/L, donc nous le couvrons
en achetant 700 contrats à terme à la date du B/L et en les revendant une fois le prix connu.~»
Corrigez-la.
\end{exercise}

\section{Problème~: Le cargo dont le prix a été fixé tardivement}

\begin{problem}[{Problème du week-end --- un cargo adossé avec deux périodes de prix}]\label{pb:m3:physical-commodity-markets:1}
Une \omterm{def:m3:physical-commodity-markets:house}{maison de négoce} achète 700\,000 barils FOB, valorisés en 2-1-2 autour de la date du B/L (jour
10) sur la base du Dated plus \qty{0.20}{\$/\barrel}, et revend le même cargo CIF Rotterdam,
valorisé en 2-1-2 autour de la date de déchargement (jour 30) sur la base du Dated plus
\qty{1.10}{\$/\barrel}. Elle paie un fret de \qty{0.65}{\$/\barrel} et une assurance de
\qty{0.03}{\$/\barrel}. Le Dated varie d'environ \qty{1.80}{\$/\barrel} par jour ouvré.

\medskip\noindent\textbf{Partie I --- L'exposition.}
\begin{enumerate}
\item Quelle est l'exposition en prix plat du trader avant le jour 8~?
\item Comment évolue-t-elle les jours 8 à 12~?
\item Quelle est-elle entre les jours 13 et 27, et pendant combien de jours atteint-elle son maximum~?
\item Comment évolue-t-elle les jours 28 à 32~?
\item Esquisser l'exposition en fonction du temps.
\end{enumerate}

\medskip\noindent\textbf{Partie II --- La couverture.}
\begin{enumerate}[resume]
\item Quels contrats à terme le trader négocie-t-il chacun des jours 8 à 12~?
\item Et les jours 28 à 32~?
\item Quelle est la position nette en contrats à terme le jour 20~?
\item Pourquoi aucune position en contrats à terme n'est-elle nécessaire avant le jour 8~?
\item Quel instrument utiliseriez-vous si les contrats à terme ne se réglaient pas sur le Dated~?
\end{enumerate}

\medskip\noindent\textbf{Partie III --- L'argent.}
\begin{enumerate}[resume]
\item Donner la marge par baril verrouillée par la couverture.
\item Donner-la pour le cargo, en dollars.
\item Estimer l'écart-type du résultat non couvert.
\item Quels risques subsistent après la couverture~?
\item Comment un retard de cinq jours dans le déchargement modifie-t-il l'exposition et la couverture~?
\end{enumerate}

\medskip\noindent\textbf{Partie IV --- Jugement.}
\begin{enumerate}[resume]
\item Pourquoi le trader pourrait-il choisir une vente CIF plutôt qu'une vente FOB~?
\item Que doit financer la ligne de crédit du trader~?
\item Pourquoi les \omterm{def:m3:physical-commodity-markets:house}{maisons de négoce} appellent-elles cette activité «~spread~» trading~?
\item Énoncer le \emph{résultat nommé}~: la marge verrouillée du cargo face à son risque non couvert.
\item En une phrase~: que gagne un trader physique~?
\end{enumerate}
\end{problem}

\medskip
\begin{omfigure}
\begin{tikzpicture}
\begin{axis}[width=11cm,height=4.6cm,
  xlabel={jour ouvré},ylabel={milliers de barils / lots},
  tick label style={font=\small},label style={font=\small},
  xmin=0,xmax=35,ymin=-800,ymax=800,grid=major,grid style={gray!25},
  legend columns=2,legend style={font=\footnotesize,at={(0.5,-0.33)},anchor=north,draw=none,fill=none,/tikz/every even column/.append style={column sep=8pt}},
  table/col sep=comma]
\addplot[omDef,very thick,const plot] table[x=day,y=exposure]{figdata/markets-3/01-physical-commodity-markets/exposure.csv};
\addplot[omThm,very thick,dashed,const plot] table[x=day,y=hedge]{figdata/markets-3/01-physical-commodity-markets/exposure.csv};
\legend{exposition physique (milliers de barils),contrats à terme détenus (lots de 1\,000 barils)}
\end{axis}
\end{tikzpicture}

\omcaption{Le cargo adossé du problème du week-end~: l'exposition au prix plat se constitue
à mesure que les prix d'achat sont fixés (jours 8--12) et se défait à mesure que les prix de vente
le sont (jours 28--32)~; la position en contrats à terme en est le miroir. Données~: le
code du chapitre.}\label{fig:m3:physical-commodity-markets:exposure}
\end{omfigure}

\section{Questions d'entretien}

\begin{interviewq}[$\star$ \iqroles{trader}]\label{iq:m3:physical-commodity-markets:1}
Quelle est la différence entre FOB et CIF, et lequel des deux porte le risque de transport~?
\end{interviewq}

\begin{interviewq}[$\star$ \iqroles{trader, chercheur}]\label{iq:m3:physical-commodity-markets:2}
Pourquoi la plupart des cargaisons physiques sont-elles valorisées sur une moyenne de plusieurs jours
plutôt qu'à un prix unique~?
\end{interviewq}

\begin{interviewq}[$\star\star$ \iqroles{trader, risque}]\label{iq:m3:physical-commodity-markets:3}
Vous avez acheté une cargaison valorisée sur la moyenne du mois suivant. Quand couvrez-vous,
et quelle quantité chaque jour~?
\end{interviewq}

\begin{interviewq}[$\star\star$ \iqroles{chercheur}]\label{iq:m3:physical-commodity-markets:4}
L'évaluation d'une \omterm{def:m3:physical-commodity-markets:pra}{agence de cotation des prix} et le prix de règlement d'une bourse sont tous deux
«~le prix~». En quoi diffèrent-ils en tant que données~?
\end{interviewq}

\begin{interviewq}[$\star\star$ \iqroles{risque}]\label{iq:m3:physical-commodity-markets:5}
Une \omterm{def:m3:physical-commodity-markets:house}{maison de négoce} affirme couvrir tout son risque de prix plat. Où peut-elle encore
perdre de l'argent~?
\end{interviewq}

\begin{interviewq}[$\star\star\star$ \iqroles{développeur, trader}]\label{iq:m3:physical-commodity-markets:6}
Concevez le service qui indique à un desk physique, chaque soir, son exposition par indice de référence
et par jour de fixation, ainsi que les contrats à terme qu'il devrait détenir.
\end{interviewq}
